% §6.2 Intersection, Complement, Difference
%
% Cross-reference resolution status:
%
% DEFINED here:
%   \label{sec:other-set-ops}                   -- this section
%   \label{def:intersection-of-sets}            -- intersection definition
%   \label{prop:intersection-fpr-fnr}           -- intersection rate proposition
%   \label{def:complement-of-set}               -- complement definition
%   \label{prop:complement-swap}                -- complement rate proposition
%   \label{def:difference-of-sets}              -- difference definition
%   \label{prop:difference-fpr-fnr}             -- difference rate proposition
%
% ALREADY DEFINED (chapter 3):
%   \Cref{ch:boolean-algebra}                   -- chapter 3
%   \Cref{sec:composition-theorem}              -- chapter 3 §3.3
%   \Cref{thm:composition}                      -- chapter 3 §3.3
%   \Cref{sec:or-not}                           -- chapter 3 §3.2
%
% ALREADY DEFINED (chapter 5):
%   \Cref{ch:approximate-sets}                  -- chapter 5
%   \Cref{def:bloom-filter}                     -- chapter 5 §5.2
%   \Cref{sec:one-sided-two-sided}              -- chapter 5 §5.4
%   \Cref{def:one-sided}                        -- chapter 5 §5.4
%
% ALREADY DEFINED (chapter 6, earlier sections):
%   \Cref{sec:union-of-sets}                    -- §6.1
%   \Cref{prop:union-fpr-fnr}                   -- §6.1
%   \Cref{eq:union-fpr}                         -- §6.1
%
% FORWARD REFS within chapter 6 (resolve in later Plan 6 tasks):
%   \Cref{sec:cancellation}                     -- §6.3 (Task 4)
%   \Cref{sec:perfect-filter-line}              -- §6.4 (Task 5)

\section{Intersection, Complement, Difference}
\label{sec:other-set-ops}

The union of \Cref{sec:union-of-sets} was the first compositional step.
The Boolean algebra of classical sets has three other operations to
account for: intersection, complement, and difference. Each lifts to
approximate sets in the same way the union did. The membership
predicate of the composed set is a Boolean combination of the operand
predicates, and the rates of the composed set follow from the
corresponding gate analysis of \Cref{ch:boolean-algebra} together with
the cross-set independence assumption already invoked in §6.1. This
section takes the three operations in turn, derives the rate formulas,
collects them in a worked numerical example, and signposts the
cancellation phenomena that §6.3 will scrutinize.

\subsection*{Intersection}

A classical intersection $A \cap B$ has the AND of the two operand
membership predicates: $\SetIndicator{A \cap B}(x) = \SetIndicator{A}(x)
\land \SetIndicator{B}(x)$. The predicate framing carries over.

\begin{definition}[Intersection of approximate sets]
\label{def:intersection-of-sets}
Let $\ASet{A}, \ASet{B} : \Universe \to \Bernoulli{\BoolType}$ be two
approximate sets over the same universe $\Universe$. Their
\emph{intersection} $\ASet{A} \cap \ASet{B}$ is the approximate set
whose membership predicate is the Boolean AND of the two operand
predicates: for every $x \in \Universe$,
\[
  (\ASet{A} \cap \ASet{B})(x) \;=\; \ASet{A}(x) \,\land\, \ASet{B}(x).
\]
The intersection approximates the objective set $A \cap B$, where $A$
and $B$ are the objective sets that $\ASet{A}$ and $\ASet{B}$
represent.
\end{definition}

The operational realization for two Bloom filters with the same
parameters and the same hash family is the bit-wise AND of the two
bit arrays, in the same spirit as the bit-wise OR for union. The
bit-AND construction has a subtlety we will flag at the end of this
subsection; the predicate-level definition above is unambiguous and is
the one the rate algebra operates on.

\begin{proposition}[Intersection rate formulas]
\label{prop:intersection-fpr-fnr}
Let $\ASet{A}$ and $\ASet{B}$ be approximate sets of $A, B \subseteq
\Universe$ with per-block rates $(\fprate_A, \fnrate_A)$ and
$(\fprate_B, \fnrate_B)$ respectively. Assume the error events of the
two operands at every query point $x$ are mutually independent (the
same cross-set independence assumption as~\Cref{prop:union-fpr-fnr}).
Then the intersection $\ASet{A} \cap \ASet{B}$ has false-negative rate
\begin{equation}
\label{eq:intersection-fnr}
  \fnrate_\cap \;=\; 1 - (1 - \fnrate_A)(1 - \fnrate_B),
\end{equation}
and false-positive rate
\begin{equation}
\label{eq:intersection-fpr}
  \fprate_\cap \;=\; \fprate_A \cdot \fprate_B
\end{equation}
on negatives at which the two operand sets agree (the bulk regime;
see below).
\end{proposition}

\paragraph{FNR derivation.}
Fix a query point $x \in A \cap B$. Both operand predicates should
report $\top$. A false negative of the intersection occurs when at
least one of the two reports $\bot$, that is, when the AND fails to be
$\top$. We compute the false-negative event as the complement of the
joint correct-report event:
\begin{align*}
  \Prob{(\ASet{A} \cap \ASet{B})(x) = \bot \Given x \in A \cap B}
  &= 1 - \Prob{\ASet{A}(x) = \top,\, \ASet{B}(x) = \top
       \Given x \in A,\, x \in B} \\
  &= 1 - \Prob{\ASet{A}(x) = \top \Given x \in A}\,
         \Prob{\ASet{B}(x) = \top \Given x \in B} \\
  &= 1 - (1 - \fnrate_A)(1 - \fnrate_B).
\end{align*}
The probability is the same for every $x \in A \cap B$, so the
intersection's false-negative rate is~\eqref{eq:intersection-fnr}. The
formula is the dual of the union FPR~\eqref{eq:union-fpr} with
$\fnrate$ in place of $\fprate$, exactly as De Morgan duality leads us
to expect.

\paragraph{FPR derivation, simplified.}
A false positive of the intersection requires both operand predicates
to report $\top$ at an $x \notin A \cap B$. For $x \notin A$ and
$x \notin B$ simultaneously, the two reports are false positives of the
operands with conditional probabilities $\fprate_A$ and $\fprate_B$, so
by cross-set independence the joint event has probability
$\fprate_A \cdot \fprate_B$, giving~\eqref{eq:intersection-fpr}. The
full intersection FPR mixes this contribution with two boundary regions
where $x$ lies in exactly one operand set; the resulting weighted
average is bounded above by $\fprate_A + \fprate_B$ and below by
$\fprate_A \cdot \fprate_B$, and the simplified form is exact when the
two operand sets coincide on the negatives of the intersection. We
adopt~\eqref{eq:intersection-fpr} as the working formula for the
intersection FPR; the full mixture appears in the
\Cref{prop:union-fpr-fnr}-style derivation that the
companion composition paper records.

\paragraph{The multiplicative shape matters.}
The intersection FPR is the \emph{product} of the operand FPRs, not a
sum. For small operand rates the contrast with the union is sharp:
two filters at $\fprate \approx 0.01$ give a union FPR near $0.02$ but
an intersection FPR near $0.0001$. Intersection drives false positives
exponentially toward zero with the number of independent filters; union
inflates them roughly linearly. The dual fact for false negatives is
that intersection's FNR adds in the same inclusion-exclusion sense the
union's FPR did, while union's FNR is multiplicative across the
regions where both operands could miss.

\subsection*{Complement}

Complement is the simplest of the three operations to analyze.

\begin{definition}[Complement of an approximate set]
\label{def:complement-of-set}
Let $\ASet{A} : \Universe \to \Bernoulli{\BoolType}$ be an approximate
set of $A \subseteq \Universe$. Its \emph{complement}
$\complement \ASet{A}$ is the approximate set whose membership
predicate is the Boolean NOT of $\ASet{A}$'s predicate:
$(\complement \ASet{A})(x) = \lnot \ASet{A}(x)$.
The complement approximates the classical complement
$\complement A = \Universe \setminus A$.
\end{definition}

\begin{proposition}[Complement-swap of rates]
\label{prop:complement-swap}
If $\ASet{A}$ has per-block rates $(\fprate_A, \fnrate_A)$, then its
complement $\complement \ASet{A}$ has rates
\begin{equation}
\label{eq:complement-swap}
  \fprate_\complement = \fnrate_A, \qquad \fnrate_\complement = \fprate_A.
\end{equation}
\end{proposition}

The result is the set-level reading of the NOT-under-noise calculation
already done in \Cref{sec:or-not} at the per-bit level, lifted from
the noisy bit to the noisy predicate. We give the two-line proof.

A false positive of $\complement \ASet{A}$ is the event that the truth
says negative and the report says positive on the complement side:
$x \in \complement A$ is false (so $x \in A$) while
$\complement \ASet{A}(x) = \top$ (so $\ASet{A}(x) = \bot$). The
combined event is $\ASet{A}(x) = \bot$ given $x \in A$, which is a
false negative of $\ASet{A}$ at rate $\fnrate_A$. So
$\fprate_\complement = \fnrate_A$. Dually, a false negative of
$\complement \ASet{A}$ is the event $x \in \complement A$ true (so
$x \notin A$) while $\complement \ASet{A}(x) = \bot$ (so
$\ASet{A}(x) = \top$). The combined event is $\ASet{A}(x) = \top$
given $x \notin A$, which is a false positive of $\ASet{A}$ at rate
$\fprate_A$. So $\fnrate_\complement = \fprate_A$. The two rates
swap, as~\eqref{eq:complement-swap} records.

The same result was proven at the noisy-bit level in \Cref{sec:or-not}.
The lifting to the set level uses the predicate framing of
\Cref{ch:approximate-sets} and nothing else. NOT is its own inverse:
applying complement twice returns the original rates,
$(\fprate_A, \fnrate_A) \mapsto (\fnrate_A, \fprate_A) \mapsto
(\fprate_A, \fnrate_A)$. The complement operation is an involution on
the rate plane $[0, 1]^2$, exchanging the two coordinate axes.

\subsection*{Difference}

Set difference is intersection with the complemented second operand.

\begin{definition}[Difference of approximate sets]
\label{def:difference-of-sets}
For approximate sets $\ASet{A}, \ASet{B}$ over $\Universe$, the
\emph{difference} $\ASet{A} \setminus \ASet{B}$ is the approximate set
whose membership predicate is $\ASet{A}(x) \land \lnot \ASet{B}(x)$;
equivalently, $\ASet{A} \setminus \ASet{B} = \ASet{A} \cap \complement
\ASet{B}$. It approximates the classical difference
$A \setminus B = A \cap \complement B$.
\end{definition}

\begin{proposition}[Difference rate formulas]
\label{prop:difference-fpr-fnr}
Under the same cross-set independence as
\Cref{prop:intersection-fpr-fnr}, the difference $\ASet{A} \setminus
\ASet{B}$ has false-positive rate
\begin{equation}
\label{eq:difference-fpr}
  \fprate_\setminus \;=\; \fprate_A \cdot \fnrate_B
\end{equation}
on the bulk negative region, and false-negative rate
\begin{equation}
\label{eq:difference-fnr}
  \fnrate_\setminus \;=\; 1 - (1 - \fnrate_A)(1 - \fprate_B).
\end{equation}
\end{proposition}

The proof is compositional. The complement of $\ASet{B}$ has rates
$(\fprate_B^\complement, \fnrate_B^\complement) = (\fnrate_B,
\fprate_B)$ by \Cref{prop:complement-swap}. Substitute these into the
intersection formulas of \Cref{prop:intersection-fpr-fnr} applied to
$\ASet{A}$ and $\complement \ASet{B}$:
\[
  \fprate_\setminus
  \;=\; \fprate_A \cdot \fprate_B^\complement
  \;=\; \fprate_A \cdot \fnrate_B,
\]
which is~\eqref{eq:difference-fpr}, and
\[
  \fnrate_\setminus
  \;=\; 1 - (1 - \fnrate_A)(1 - \fnrate_B^\complement)
  \;=\; 1 - (1 - \fnrate_A)(1 - \fprate_B),
\]
which is~\eqref{eq:difference-fnr}.

The formulas mix the two operands' rates in a characteristic way:
the difference FPR pairs $\fprate_A$ with $\fnrate_B$ (the swapped
complement-side FPR), and the difference FNR pairs $\fnrate_A$ with
$\fprate_B$ (the swapped complement-side FNR). Both operands'
\emph{both} rates appear in the answer; even an FNR-zero $\ASet{A}$
inherits false negatives from the difference when $\ASet{B}$ has a
nonzero FPR, because elements of $A \setminus B$ that are false
positives of $\ASet{B}$ are wrongly excluded.

\subsection*{Worked numerical example}

Let $\ASet{A}$ and $\ASet{B}$ be two Bloom filters over the same
universe with the same parameters and operand rates
\[
  \fprate_A = 0.01, \qquad \fnrate_A = 0,
  \qquad
  \fprate_B = 0.02, \qquad \fnrate_B = 0.
\]
Both operands are FNR-zero (\Cref{def:one-sided}), so they sit on the
FNR-zero axis of the perfect filter line. The intersection has
\[
  \fprate_\cap \;=\; \fprate_A \cdot \fprate_B
                \;=\; 0.01 \times 0.02
                \;=\; 0.0002,
  \qquad
  \fnrate_\cap \;=\; 1 - (1 - 0)(1 - 0) \;=\; 0.
\]
The intersection of two FNR-zero Bloom filters is itself FNR-zero, and
its false-positive rate is smaller than either operand's by an order of
magnitude (an absolute drop of nearly a factor of $50$ relative to the
worse operand). The contrast with the union's $\fprate_\cup = 0.0298$
of \Cref{sec:union-of-sets} is striking: composing two filters
multiplicatively reduces the false-positive rate, while composing them
disjunctively inflates it. The two operations are exact duals of one
another, and the choice of which to use is dictated by what one wants
the composed filter to assert.

The complement of $\ASet{A}$ has rates
$(\fprate_\complement, \fnrate_\complement) = (\fnrate_A, \fprate_A)
= (0, 0.01)$. The complement has moved $\ASet{A}$ off the FNR-zero
axis and onto the FPR-zero axis: what was a one-sided false-positive
machine is now a one-sided false-negative machine. The next section
returns to this geometric observation.

The difference $\ASet{A} \setminus \ASet{B}$ has
\[
  \fprate_\setminus \;=\; \fprate_A \cdot \fnrate_B \;=\; 0.01 \times 0 \;=\; 0,
  \qquad
  \fnrate_\setminus \;=\; 1 - (1 - 0)(1 - 0.02) \;=\; 0.02.
\]
The difference of two FNR-zero filters has dropped its false-positive
rate to zero (because $\fnrate_B = 0$ kills the FPR product) and picked
up a false-negative rate equal to $\fprate_B$. The structural reason
is that $\ASet{B}$'s false positives become misclassified exclusions in
$\ASet{A} \setminus \ASet{B}$: an element of $A \setminus B$ that
$\ASet{B}$ erroneously reports as a member of $B$ will be filtered out
by the difference even though it should be kept. The difference
exemplifies the cancellation phenomena §6.3 will scrutinize.

\subsection*{A caveat on bit-AND of Bloom filters}

The intersection of \Cref{def:intersection-of-sets} is a predicate-level
operation: query $\ASet{A}$, query $\ASet{B}$, AND the two Boolean
outputs. For two Bloom filters with the same parameters and the same
hash family, this predicate intersection is the canonical definition
and the rate formulas of \Cref{prop:intersection-fpr-fnr} apply to it.
The operational shortcut of bit-wise AND of the two underlying bit
arrays is a different construction. The bit-AND of two Bloom filters
gives a bit array whose query behaviour reports $\top$ when the AND of
the bit arrays has $1$ at every queried position; this is the same as
both filters individually reporting $\top$, so on items inserted into
\emph{both} sets the bit-AND query matches the predicate AND. On items
inserted into one set only, however, the bit-AND filter can still
report $\top$ when both bits happen to be set in the AND'd array (one
from the legitimate insertion, one from a coincidental hash collision
of an insertion of the other set), and the resulting filter is not in
general a Bloom filter for $A \cap B$. The bit-AND is an operational
shortcut that coincides with the predicate intersection in some regimes
and not others; the chapter notebook explores the discrepancy
empirically, and we will not rely on bit-AND in the rate analyses to
come.

\subsection*{What comes next}

The four operations of \S6.1 and \S6.2 (union, intersection,
complement, difference) suffice to express every Boolean combination of
approximate sets, with closed-form rate formulas derived from the
composition theorem of \Cref{thm:composition} applied to the
corresponding Boolean gate. \Cref{sec:cancellation} takes up the
phenomena that the difference example already foreshadowed: composition
can move an operand off the perfect filter line, can swap which side of
the line a one-sided set sits on, and can introduce both rates as
nonzero even when each input had one rate set to zero by construction.
\Cref{sec:perfect-filter-line} then restates the line itself as a
formal subalgebra: the FNR-zero axis is closed under union and
intersection, the FPR-zero axis is closed dually, and the two axes
together form a non-convex set whose closure properties this chapter's
operations either preserve or break.
